\begin{afterword}
\summaryhead

The author recalls a synagogue's annual appeal, where members called out their pledges, and sets beside it an online appeal the author ran for a nonprofit, sent to an audience of atheists, libertarians, technophiles and programmers. Critics posted their reasons for not giving; donors gave but said nothing on the lists, so each donor saw only objections.

Why can a flying-saucer cult gather followers when ``rationalists'' cannot? Cults are held together by pluralistic ignorance, evaporative cooling and affective death spirals. Rationalists should not copy these, but a true rationalist ought to be effective, and knowing more should not make a group do worse. The author diagnoses half-rationality: a culture that practises and rewards only disagreement, as in a conference where praise at the microphone would be horrifying, and a shame of strong feeling. The remedy is to train agreement as well as dissent, to feel emotions that fit the facts, and to voice support publicly, as the synagogue's members did.

\responsehead

The question is fair and the remedy has evidence behind it. The diagnosis has the author's memory behind it.

The evidence is one fundraising appeal, told years later by the person who ran it, and the author's impressions of the community. From these the post extracts ``the two major forces'' that keep the ``atheist\slash{}libertarian\slash{}technophile'' crowd from coordinating. The title says the crowd ``Can't''; the body, more cautious, says ``Perhaps that is why.'' Much of the rest reassembles earlier posts: pluralistic ignorance, evaporative cooling, ``Knowing About Biases Can Hurt People'' and ``Feeling Rational.'' The outside record is mixed. In the year before the post, supporters of Ron Paul, a former Libertarian Party nominee, raised over \$4.2 million online in one day, and a British atheist bus campaign that failed on its first try raised over \pounds100,000 in four days on its second. Such crowds can coordinate, at least sometimes. Whether it does so less than others, the post does not measure.

Two findings favour the post. Its remedy, public support voiced alongside objections, has experimental backing: in a field experiment on charitable giving, telling new donors what others had given raised their gifts. And the Raelians, the post's example of a group held together by nonsense, include many who call their creed an ``atheistic religion.'' So the contrast is not between atheists and believers, which suits the post's thesis about culture better than it suits the post's framing.

The paragraph saying the ``Light Side will always be divided and weak'' and that ``the future inevitably belongs to the Dark'' states a view the post goes on to reject; a quick reader should not take it for the conclusion. The objection held over at the start, that the critics may have had a point, is answered at the end with a rule about timing: object at another time, or beside public support.

In short: a diagnosis of the author's community drawn from one fundraiser and the author's impressions, under a title firmer than its hedges, with a remedy, public support beside public objections, that an experiment on giving supports.
\end{afterword}
